\chapter{Project files}\label{ch:files}

A project is one \file{.psn} file: UTF-8 JSON with line feeds, readable in any text
editor and easy to compare in a version-control diff. It holds the network, its
variables, the drawing and the study cases. Results are not saved; run the cases again
after opening a project.

\section{Structure}

\begin{reftable}{@{}>{\raggedright\arraybackslash}p{2.6cm}X@{}}
\toprule
\sffamily\bfseries Key & \sffamily\bfseries Contents \\ \midrule
\code{format}, \code{version} & \code{"psa-project"} and the format version, 2 today. \\
\code{meta} & Title and notes of the project. \\
\code{network} & \code{s\_base\_mva}, \code{f\_hz}, the base reference
  (\code{base\_ref}: \code{bus} and \code{u\_base\_kv}), the \code{buses} and the
  \code{components}. Each component is an object with its \code{type}, its \code{id} and
  every field by its name. Complex numbers are written \code{[re, im]}. \\
\code{parameters} & \code{script\_path}, the \code{project\_vars} as
  \code{[name, expression]} pairs in order, and the \code{bindings}
  (\code{"LD4.p\_mw": "158.5 + 0.1*A"}). \\
\code{layout} & The \code{grid} step, the \code{items} (each element's position,
  rotation, bus orientation and length, and a dragged label offset) and the \code{wires}
  (the bends of each wire, keyed \code{"COMPONENT:terminal"}). \\
\code{cases} & The study cases: \code{id}, \code{study}, \code{title}, \code{options},
  \code{overrides} and \code{exclude}. Complex options are written
  \code{\{"re": ..., "im": ...\}}. \\
\code{view} & The display state, such as the annotations toggle. \\
\bottomrule
\end{reftable}

A case as it appears in \file{examples/lf\_220kv.psn}:

\begin{terminal}
{"id": "LFd", "study": "loadflow",
 "title": "LFd) Load flow with P_LD4 increased by 10 %",
 "options": {}, "overrides": {"LD4.p_mw": "1.1*value"}}
\end{terminal}

\section{Compatibility}

\begin{itemize}
  \item \textbf{Older files} are upgraded when opened. Version~1 kept per-study options
    instead of cases, and those become cases. The file is written in the current format
    the next time you save it.
  \item \textbf{New fields} have defaults, so a file written before a field existed opens
    with the default. A field the program does not know is ignored.
  \item \textbf{Only the network is required.} A file without a \code{layout} opens with
    an automatic layout: buses from left to right by distance from the base reference
    bus, branches between their buses and shunt elements below theirs.
\end{itemize}

\section{The variable script}

When a project links a variable script, the script stays a separate \file{.py} file. The
project keeps its path and, under \code{project\_vars} and in the fields, the values it
gave last time. Opening the project never runs the script without asking
(\cref{sec:script}).

\section{Saving and examples}

\keys{\ctrl + S} saves to the project's file. An untitled project, or one opened as a
copy such as \cli{run\_gui.py --example}, asks for a name instead, so an example is never
overwritten by accident. The window title shows the file's name, with an asterisk while
there are unsaved changes. Closing the window, opening another project or starting a new
one asks first whether to save them.
